\section{The raw cancellation invariant}
\label{sec:raw}

This section records the exact algebraic prerequisite from the preceding
positive-cancellation work \cite{IncomingPositive}.  The use of unsigned
literal occurrences is essential.  These counts record intersections with
specified belt discs; they are not exponent sums of freely reduced words.

\begin{definition}[Positive singleton pivot and stopping set]
Let $A=(a_{rc})$ be a finite nonnegative integer matrix with rows $R$ and
columns $C$.  A positive singleton pivot chooses $(r,c)$ with $a_{rc}=1$,
deletes row $r$ and column $c$, and changes each surviving entry to
\[
 a'_{id}=a_{id}+a_{ic}a_{rd}\qquad(i\ne r,\ d\ne c).
\]
A nonempty set $S\subseteq R$ is a \emph{stopping set} if
\[
 s_c:=\sum_{r\in S}a_{rc}\in\{0\}\cup\{2,3,\ldots\}
 \quad\text{for every }c\in C.
\]
A \emph{column leaf} has total live column occurrence one.  Pivoting on
it creates no fill in the surviving rows and is called a peeling move.
\end{definition}

The formula deliberately has a plus sign.  It is the conservative raw
intersection update: algebraic cancellation of opposite letters is not
part of this model.  The theorem below cannot be extended to arbitrary
subtractive elimination merely by reusing its proof.

\begin{lemma}[A stopping set survives positive pivots]
\label{lem:stopping-invariant}
If $A$ contains a stopping set, then every legal positive singleton pivot
leaves a nonempty stopping set on surviving rows.
\end{lemma}
\begin{proof}
Let $(r,c)$ be the pivot and let $S$ be a stopping set.  If $r\notin S$,
retain $S$.  For each surviving column $d$ its new total is
$s_d+s_c a_{rd}$.  If this is positive, either $s_d\ge2$ or
$s_c a_{rd}\ge2$, so it is not one.

If $r\in S$, the pivot entry is one and $s_c\ne1$, hence $s_c\ge2$.
Another row of $S$ meets $c$, so $S'=S\setminus\{r\}$ is nonempty.  Its
new total in column $d\ne c$ is
\[
 \sum_{i\in S'}a'_{id}
 =s_d-a_{rd}+(s_c-1)a_{rd}
 =s_d+(s_c-2)a_{rd}.
\]
If $s_d=0$, then $a_{rd}=0$ and the total is zero.  Otherwise $s_d\ge2$
and the total is at least two.  Thus $S'$ is a stopping set.
\end{proof}

\begin{theorem}[Complete raw cancellation equals peeling; prior result]
\label{thm:raw-completeness}
For every finite nonnegative integer matrix, the following are equivalent:
there exists a complete positive singleton sequence; column-leaf peeling
deletes every row; and there is no nonempty stopping set.  Peeling's
terminal row set is independent of queue order and is the greatest
stopping set, with the empty set allowed as the greatest residual.
\end{theorem}
\begin{proof}
A nonempty stopping set excludes a complete positive sequence by
Lemma~\ref{lem:stopping-invariant}.  If peeling stops with live rows, their
column totals are zero or at least two, so these rows form a stopping set.
Consequently the absence of stopping sets forces peeling to finish, and
peeling itself is a positive singleton sequence.

Any stopping set in the original matrix survives every peeling move
unchanged: a live column of total one cannot meet it.  Hence every original
stopping set is contained in the terminal row set.  That terminal set is
itself a stopping set when nonempty.  It is therefore the greatest one and
does not depend on the order of valid leaves.  Equivalently, unions of
stopping sets are stopping sets: a column meeting the union meets one of
its stopping subsets, which already contributes at least two.
\end{proof}

\begin{corollary}[Heredity under row deletion]
\label{cor:row-heredity}
Every row restriction of a completely peelable matrix is completely
peelable.
\end{corollary}
\begin{proof}
A stopping subset in the restricted matrix would be exactly the same
stopping subset in the original matrix.  Apply
Theorem~\ref{thm:raw-completeness}.
\end{proof}

Heredity is why a successful independent set of dual deletions can later
be extended to a spanning tree.  It does not say that an arbitrary primal
column deletion preserves success; deleting columns can remove the only
available leaves.  The new geometric theorem repairs precisely that
asymmetry.

\subsection{Saturation and immutable-source replay}

The map $x\mapsto\min\{2,x\}$ preserves whether a nonnegative sum or
product is zero, one, or at least two.  Therefore all pivot-admissibility
and stopping-set decisions can use the three values $0,1,2^+$ without
expanding large counts.  In contrast, this saturation does not preserve
parity: three becomes two.  Geometric parity conditions must be proved on
the uncapped source occurrences first.

During peeling, source coefficients never change.  Keep, for each column,
the number of remaining incident rows and the exclusive-or of their
identifiers.  If there is one owner and its source coefficient is one,
the column is a leaf.  Deleting that row updates each incident column
once.  With indexed adjacency lists this is
$O(|R|+|C|+I)$ work and storage, where $I$ is the number of nonzero
incidences.  Dictionary lookup, label encoding, and graph construction
costs are separate implementation considerations.

A positive certificate lists distinct row/column pivot pairs.  A checker
can validate them directly against immutable source rows, or use the
existing timestamp criterion: at the deletion time of row $r$, every
other owner of its chosen column must already be gone, and the coefficient
at $r$ must be one.  Successful production does not authorize a checker
to trust an unauthenticated occurrence matrix.

\subsection{Why signed homology is insufficient}

A curve traversing a belt disc three times can have algebraic intersection
one.  Its raw count is still three, so it is not a geometric singleton in
the current charts.  Likewise, a determinant equal to one can certify a
homological condition without exhibiting the disjoint attaching curves
needed for handle cancellation.  The incoming homological progress work
can guide surface choices, but its output cannot be silently substituted
for the raw matrix in this theorem.  Later sections discuss a different,
carefully scoped terminal use of verified group presentations.
